\section{The all-higher-level contact identity}
\label{sec:contacts}

The key estimate is obtained by keeping the anchor fixed while the
higher return level varies. Write
\[
 m_d=(P^{\circ p^d})'(\beta),\qquad A_d=v(m_d-1),
 \qquad \beta\in\Pi_d.
\]
We first prove that $m_d-1$ is nonzero, before using it as a
denominator.

\begin{proposition}[Multiplier and average contacts]\label{prop:contacts}
For every $d\ge1$ and $\beta\in\Pi_d$,
\begin{equation}\label{eq:multiplier-sum}
 A_d=\delta_{d-1}+
       \sum_{j=0}^{d-1}(p-1)p^{d-j-1}\delta_j<\infty.
\end{equation}
The value $A_d$ is independent of $\beta$. For every $e>d$,
\begin{equation}\label{eq:contact-equality}
 \frac1{p^e}\sum_{\alpha\in\Pi_e}v(\beta-\alpha)
 =c_d:=\frac{p-1}{p^{d+1}}A_d
 \ge b_d:=d r^3+r^2\left(1+\frac1p\right).
\end{equation}
\end{proposition}
\begin{proof}
Differentiate the factorization of Lemma~\ref{lem:small-factor}
at level $d$ and evaluate at $\beta$. Since $M_d(\beta)=0$,
\begin{equation}\label{eq:multiplier-product}
 m_d-1=F_{d-1}(\beta)M_d'(\beta)V_d(\beta).
\end{equation}
The first factor on the right is nonzero by the least-period
condition, the second by the simple root product for $M_d$, and
the third because it is a unit. Hence $m_d\ne1$. Taking valuations
and using Lemma~\ref{lem:index} gives
\begin{align*}
 A_d
 &=v(F_{d-1}(\beta))+
      \sum_{\gamma\in\Pi_d\setminus\{\beta\}}v(\beta-\gamma)\\
 &=\delta_{d-1}+\sum_{j=0}^{d-1}(p-1)p^{d-j-1}\delta_j.
\end{align*}
The chain-rule expression for $m_d$ is the product of $P'$ over
all points of the cycle. A different anchor only permutes its
factors, proving independence of $A_d$.

Now let $e>d$. Both $F_e$ and $F_{e-1}$ vanish at $\beta$.
Differentiating $F_e=F_{e-1}Q_e$ there gives
\begin{equation}\label{eq:derivative-quotient}
 F_e'(\beta)=F_{e-1}'(\beta)Q_e(\beta).
\end{equation}
Because $P^{\circ p^d}$ fixes $\beta$, the chain rule and
characteristic $p$ yield
\[
 F_e'(\beta)=m_d^{p^{e-d}}-1=(m_d-1)^{p^{e-d}},\qquad
 F_{e-1}'(\beta)=(m_d-1)^{p^{e-d-1}}\ne0.
\]
The last exponent is one when $e=d+1$, so that boundary case also
allows division in \eqref{eq:derivative-quotient}. We obtain
\begin{equation}\label{eq:higher-quotient}
 Q_e(\beta)=(m_d-1)^{(p-1)p^{e-d-1}}.
\end{equation}
Distinct least periods make $\Pi_d$ and $\Pi_e$ disjoint. Using
$Q_e=M_eV_e$, the unit value of $V_e(\beta)$ and the exact root
product for $M_e$, equation~\eqref{eq:higher-quotient} becomes
\begin{equation}\label{eq:contact-total}
 \sum_{\alpha\in\Pi_e}v(\beta-\alpha)
 =(p-1)p^{e-d-1}A_d.
\end{equation}
All summands are finite. Division by the real number $p^e$ proves
the equality in \eqref{eq:contact-equality}.

Finally Lemma~\ref{lem:displacement} and
\eqref{eq:multiplier-sum} give
\begin{align*}
 \frac{A_d}{p^d}
 &\ge \frac{r(p^{d-1}+1)}{p^d}
     +\frac r{p^d}\sum_{j=0}^{d-1}
        (p-1)p^{d-j-1}(p^j+1)\\
 &=r\left(1+\frac1p+d\frac{p-1}{p}\right).
\end{align*}
Multiplication by $(p-1)/p=r$ proves the lower bound in
\eqref{eq:contact-equality}. In particular $c_d\to+\infty$.
\end{proof}

\begin{corollary}[Aligned native-orbit coupling]\label{cor:coupling}
For every $e>d$, there is a coupling of $\mu_e$ and $\mu_d$ all
of whose support pairs have distance at most $|s|^{c_d}$.
Moreover $d_H(\Pi_e,\Pi_d)\le|s|^{c_d}$.
\end{corollary}
\begin{proof}
Fix $\beta\in\Pi_d$. One contact in
\eqref{eq:contact-equality} is at least its real average. Choose
$\alpha\in\Pi_e$ realizing such a contact. Lemma~\ref{lem:isometry}
gives
\[
 |P^{\circ i}(\alpha)-P^{\circ i}(\beta)|
 =|\alpha-\beta|\le |s|^{c_d}\qquad(i\ge0).
\]
The real probability
\begin{equation}\label{eq:aligned-coupling}
 \pi_{e,d}=p^{-e}\sum_{i=0}^{p^e-1}
    \delta_{(P^{\circ i}(\alpha),P^{\circ i}(\beta))}
\end{equation}
has first marginal $\mu_e$. Each point of $\Pi_d$ appears
$p^{e-d}$ times in the second coordinate, giving marginal $\mu_d$.
The same matched iterates cover both cycle sets and prove the
Hausdorff bound. No compatible phases across different pairs of
levels are required.
\end{proof}
